% Esquisse d'un programme (cote 119, copy 2): archive cote 119, site page 130 (PDF page 131), typescript page 52.
% typescript. Draft compared: esquisse/pages/esq_119_054.tex (copy 1, site 54).
% AI-transcribed by Claude Opus 5.5 (claude-opus-5-5, adaptive thinking, effort high), Message Batches API, batch esq119_full (msgbatch_015xqDHkTumDysRF1qngZhRJ), 2026-11-13; not yet checked by eye.
\scanpage{119/130}
\mspage{52}

(\textsuperscript{1}) L'expression ``hors de portée'' ici (et encore plus loin pour
une question toute différente), que j'ai laissé passer en allant
à l'encontre d'une réticence, me paraît décidément hâtive et sans
fondement. J'ai pu constater déjà en d'autres occasions que lors-%
que des augures (ici moi-même !) déclarent d'un air entendu
(ou dubitatif) que tel problème est ``hors de portée'', c'est là
au fond une affirmation entièrement subjective. Elle signifie
simplement, à part le fait que le problème est censé ne pas être
résolu encore, que celui qui parle est à court d'idées sur la
question, ou de façon plus précise sans doute, qu'il est devant
elle sans sentiment ni entrain, qu'elle ``ne lui fait rien'' et
qu'il n'a aucune envie de faire quelque chose avec elle - ce qui
souvent est une raison suffisante pour vouloir en décourager au-%
trui. Cela n'a pas empêché qu'à l'instar de M. de la Palisse, et
au moment même de succomber, les belles et regrettées conjectures
de Mordell, de Tate, de Chafarévitch étaient toujours réputées
``hors de portée'', les pauvres ! - D'ailleurs, dans les jours déjà
qui ont suivi la rédaction du présent rapport, qui m'a remis en
contact avec des questions dont je m'étais quelque peu éloigné au
cours de l'année écoulée, je me suis aperçu d'une nouvelle pro-%
priété remarquable de l'action extérieure d'un groupe de Galois
absolu sur le groupe fondamental d'une courbe algébrique, qui m'avait
échappée jusqu'à présent et qui sans doute constitue pour le moins
un nouveau pas en avant vers la formulation d'une caractérisation
algébrique de $\mathrm{Gal}(\bar{\underline{Q}}/\underline{Q})$. Celle-ci, avec la ``conjecture fondamen-%
tale'' (mentionnée au par. 3 ci-dessous) apparaît à présent comme
la principale question ouverte pour les fondements d'une ``géométrie
algébrique anabélienne'', laquelle depuis quelques années, repré-%
sente (et de loin) mon plus fort centre d'intérêt en mathématiques.

(\textsuperscript{2}) Je puis faire exception pourtant d'un autre ``fait'', du temps
où, vers l'âge de douze ans, j'étais interné au camp de concentra-%
tion de Rieucros (près de Mende). C'est là que j'ai appris, par une
détenue, Maria, qui me donnait des leçons particulières bénévoles,

\margin{130}
